\subsection{The framed presentation and its real-place defect}
\label{dp:section}

Let \(O\) be the ring of integers of a finite extension \(E/\mathbf Q_\ell\),
with residue field \(k\) of odd characteristic. The two presentations used
here are presentations of \emph{framed} representations. In rank three
the form-preserving group in the two dual blocks has Lie algebra
\(\mathfrak {gl}_3\). In rank four, fixing the multiplier gives
\(\mathfrak {sp}_4\). Write \(M_j\) for the residual Lie algebra of
weight \(j\), with Galois action by the adjoint of the residual point.
The grading records the weight of an entry \(ab\) as \(w_a-w_b\).
In particular the tangent vector that reads that entry is indexed here
by the weight of the entry it reads.

Write \(G=G_{\mathbf Q,S}\) for the Galois group unramified outside the
finite primes \(S\) and infinity. Partition \(S=T\sqcup Q\), with
\(\ell=v_0\in T\), and use the full local decomposition groups \(G_v\)
with their maps to \(G\). At \(v\in T\) a complete local augmented
algebra \(B\) supplies a simultaneous family \(\rho_v\) of local
representations. In rank three \(T=\{\ell\}\) and \(B\) is the
prescribed quotient of the local ring. In rank four
\(T=\{\ell\}\cup S_b\). At \(q\in Q\) representations must kill
\(N_q=I_q\cap\ker\bar r\), where \(I_q\) is inertia. In both
applications \([I_q:N_q]\) is prime to \(\ell\). Let \(r_0\) be the
chosen integral lift. An object over a local \(O\)-algebra \(C\)
consists of \(i:B\to C\), a lift \(r\) with the prescribed form and
multiplier, and congruence-to-identity isometries \(\alpha_v\), with
\[
 \alpha_{v_0}=1,\qquad
 r(s)\alpha_v=\alpha_v\,i(\rho_v(s))\quad(s\in G_v).
\]
All matrix-valued maps are continuous. Representations over finite
Artin \(O\)-algebras use the discrete topology.

\begin{proposition}[Relative generators and relations]\label{dp:relative}
Suppose that \(B\) has a finite weighted power series presentation, the
entries of \(\rho_v\) have the indicated weights, and the base point
has weight zero. The framed ring \(U\) has a presentation over \(B\)
by \(g_j\) variables and at most \(o_j\) relations of weight \(j\).
The following formulas describe both integers. Put \(M=M_j\), write
\(dy(s)=s y-y\), and set
\[
 {\cal T}(M)=\{(z,(y_v)):\ z\in Z^1(G,M),\
 z|_{N_q}=0,\ z|_{G_v}=dy_v,\ y_{v_0}=0\}.
\]
Then \(g_j=\dim_k{\cal T}(M)\). An obstruction cocycle is a continuous
triple \((c,(b_v),(b_q))\), where \(c\in Z^2(G,M)\) and
\(db_w=c|_{G_w}\) for \(w\in S\). Divide by the triples
\[
 \bigl(dh,(h|_{G_v}-dy_v)_v,(h|_{G_q}-e_q)_q\bigr),
 \qquad e_q\in Z^1(G_q,M),\quad e_q|_{N_q}=0.
\]
The dimension of the resulting space \({\cal O}b(M)\) is \(o_j\).
The variables can be chosen among entries of \(r(g)-r_0(g)\) and
\(\alpha_v-1\). Zero equations can be appended in each weight to
give exactly \(o_j\) listed equations.
\end{proposition}

\begin{proof}
For clarity, the usual finite presentation argument starts with finitely
many matrices. The maximal pro-\(\ell\) quotient of \(\ker\bar r\) in
\(G\) is topologically finitely generated. Indeed its Frattini quotient
is dual to \(H^1(\ker\bar r,\mathbf F_\ell)\), which is finite by
finiteness of number-field cohomology with ramification in a finite
set. The kernel of this maximal quotient is characteristic in
\(\ker\bar r\). Every Artin lift factors through the corresponding
quotient of \(G\), since the congruence kernel of a matrix group over
a finite local \(O\)-algebra is an \(\ell\)-group. Choose topological
generators there, adjoining lifts of the finite residual quotient.
Start with \(B\) and matrix-entry variables shifted by their base
values for the generators and frames. Invertibility holds modulo the
maximal ideal and hence in the power series ring. The free profinite
group on the generators maps continuously to its matrix group by
passage through the finite coefficient quotients. Impose its relators,
the form, block, multiplier, inertia and frame equations. These are
closed ideal conditions and give precisely the stated universal
property, first over Artin rings and then over their inverse limits.
The noetherian ring already contains the limits of its ideals.
Every matrix-product summand has the weight of the resulting entry.
The base lift preserves the weight spaces and normalizes the form and
block conditions. Thus these equations respect the matrix weights,
including the relations obtained by profinite limits. This weight
argument works termwise modulo powers of the maximal ideal.

Hold \(B\) fixed at its residual point and evaluate on \(k[\epsilon]\),
\(\epsilon^2=0\). Write \(r=(1+\epsilon z)\bar r\) and
\(\alpha_v=1-\epsilon y_v\). Multiplication is exactly the cocycle
identity and the frame equations are the equations defining
\({\cal T}(M)\). In particular no quotient by global conjugation
occurs. Matrix-entry evaluations span the dual of this finite
dimensional space because they separate all the pairs. Choose a basis
among these evaluations in each weight. With the corresponding
variables \(t_i\), the map \(P=B[[t_1,\ldots,t_g]]\to U\) is
surjective by successively lifting modulo powers of the maximal ideal.
Its kernel \(I\) lies in \(\mathfrak m_B P+\mathfrak m_P^2\), by the
chosen basis. The kernel is closed under taking weight components:
substitution of the homogeneous images is componentwise modulo every
power, and both rings are separated. Therefore \(I/\mathfrak m_PI\)
has a homogeneous basis lifting to homogeneous generators. Finiteness
of this vector space permits only finitely many weights. Nakayama's
lemma gives generation of the full ideal.

For the bound, put \(C=I/\mathfrak m_PI\). In
\(P/\mathfrak m_PI\) the maximal ideal annihilates \(C\). Choose \(h\)
so large that \(C\) meets the \(h\)-th maximal-ideal power trivially.
This follows from Krull intersection and finite dimension of \(C\).
After truncating by that power, the extension onto \(P/(I+\mathfrak
m_P^h)\) therefore still has kernel \(C\). Lift the universal matrices
as functions across this finite extension, preserving their form,
and lift the frames. A form-preserving congruence matrix admits such
lifts: lift its entries and correct the form error to first order
by subtracting half its adjoint-symmetrization error in the kernel.
In the paired blocks one may equivalently lift the first block and
use its inverse transpose in the second. At \(v\in T\) the specified
lift is supplied by the structural map of \(B\). At \(q\in Q\)
there is a local lift killing \(N_q\): \(G_q/N_q\) is an extension
of \(\widehat{\mathbf Z}\) by the finite prime-to-\(\ell\) group
\(I_q/N_q\). Averaging on the finite subgroup and
\(\ell\)-cohomological dimension one of \(\widehat{\mathbf Z}\)
give \(H^2(G_q/N_q,M)=0\), which kills the local multiplication
error. Locally constant lifts of matrices exist since all the rings
in this step are finite.

For a set-theoretic lift \(\widetilde r\), the excess of
\(\widetilde r(s)\widetilde r(t)\widetilde r(st)^{-1}\) over 1
is a Lie-algebra-valued two-cocycle with coefficients in \(C\).
Comparison with the true local lifts, conjugated by the lifted
frames, gives the cochains \(b_w\). Changing the global lifts
adds \(dh\), changing frames adds a local coboundary, and changing
a local lift killing \(N_q\) adds a cocycle vanishing there.
Conversely these changes exhaust the displayed coboundaries. The
frame at \(v_0\) may be restored to 1 by common conjugation. Thus
the class after a linear pushout \(C\to k\) vanishes exactly when
the point lifts over \(B\) to that pushout. The component supported
on \(C_j\) has its matrix values in \(M_j\): choose homogeneous lifts
of the entries through the graded truncation and take the coefficient
of a degree-\(j\) kernel element in each error. For the auxiliary
local corrections one may project cochains onto the matching weight
since the residual action preserves every weight. Equivalently first
push out to a homogeneous kernel line, where only entries of its own
weight admit a nonzero homogeneous correction.

The resulting linear map \(C_j^\vee\to{\cal O}b(M_j)\) is injective.
Indeed a zero obstruction gives \(U\to A\) lifting the truncated
point in the pushed-out extension \(A\). Its composite with \(P\to U\)
and the canonical map \(P\to A\) differ by a \(B\)-derivation into
the kernel \(k\). This derivation kills
\(\mathfrak m_BP+\mathfrak m_P^2\), hence \(I\). On \(I\) the first
map is zero and the second map is exactly the chosen functional.
This proves the bound, including the case of a zero obstruction
space. The padding by zero relations makes no minimality assertion.
\src{src/Development/AutomorphicArithmetic/AutomorphicVarieties/CommutativeAlgebra/LemmaP.lean}{256}{The relative presentation argument}
and \src{src/Development/AutomorphicArithmetic/AutomorphicVarieties/CommutativeAlgebra/FramedTwoCoboundaries.lean}{249}{its precise functor and indexing}.
\end{proof}

\begin{proposition}[The equation count]\label{dp:count}
For a finite \(k[G]\)-module \(M\) on which every \(N_q\) acts
trivially, the above spaces satisfy
\[
 \dim{\cal O}b(M)-\dim{\cal T}(M)=
 \dim M^{c=1}+\dim H^0(G,M^\vee(1)),
\]
where \(c\) is complex conjugation and \(M^\vee=\operatorname{Hom}_k(M,k)\).
\end{proposition}

\begin{proof}
Let \(L_q\subset H^1(G_q,M)\) be the inflated classes from
\(G_q/N_q\) and let \(A^1\) be the product of \(H^1(G_v,M)\)
for \(v\in T\) and \(H^1(G_q,M)/L_q\) for \(q\in Q\).
Write \(r_1:H^1(G,M)\to A^1\) and
\(r_2:H^2(G,M)\to\prod_{w\in S}H^2(G_w,M)\).
Sending a pair in \({\cal T}\) to its global class has image
\(\ker r_1\). For any class there choose the local primitives and
subtract the global coboundary of the primitive at \(v_0\).
Over the zero class the choices are local invariant vectors at
\(T\) modulo the simultaneous global invariant vector, as is also
seen by starting with an arbitrary global primitive. Forgetting the
local cochains in \({\cal O}b\) has image \(\ker r_2\) by definition,
and kernel \(\operatorname{coker}r_1\), by subtracting a global
primitive of \(c\). Thus
\[
 \dim{\cal T}=\dim\ker r_1+\sum_{v\in T}h^0_v-h^0_G,\qquad
 \dim{\cal O}b=\dim\operatorname{coker}r_1+\dim\ker r_2.
\]
For \(q\in Q\), averaging over \(I_q/N_q\) expresses \(L_q\)
as the coinvariants of Frobenius on \(M^{I_q/N_q}\).
The kernel and cokernel of Frobenius minus 1 have the same dimension,
so \(\dim L_q=h_q^0\). The final segment of Poitou--Tate identifies
\(\dim\operatorname{coker}r_2=\dim H^0(G,M^\vee(1))\).
The real place has zero positive-degree cohomology since \(\ell\)
is odd. With \(\chi=h^0-h^1+h^2\), the remaining terms are
\(\chi_G-\sum_{w\in S}\chi_w\). The local Euler characteristics
are \(-\dim M\) at \(\ell\) and zero at all other finite primes.
The global Euler characteristic is \(-\dim M^{c=-1}\).
These formulas give the assertion. For the finite module formulas
see Milne~\cite{milneADT},
I, Theorems 2.8, 4.10 and 5.1. Passing from group orders to
\(k\)-dimensions divides all logarithms by \(\log |k|\).
\end{proof}

In rank three the auxiliary datum supplies \(g\) with
\(\eta(g)=\bar\Psi(g)=1\) and \(\bar\chi_\ell(g)=-1\).
Indeed its field-disjointness condition gives \(h\) in the joint
kernel of the residual cyclotomic and finite characters, outside
\(\ker\eta\). Hence \(\bar\Psi(h)=1\) and \(\eta(h)=-1\).
Since \(\eta(c)=-1\), \(\Psi(c)=1\), take \(g=ch\).
It acts as minus the identity on the twisted dual. On
\(\mathfrak {gl}_3(k)\) the action of \(c\) is negative transpose
across the antidiagonal. Its fixed part in weight zero is the line
\(\operatorname{diag}(a,0,-a)\). In either first off-diagonal the
two entries are negative partners, giving one line, and each outer
corner is its own negative, giving zero. The relative list therefore
has \(g_j+{\bf1}_{|j|\le1}\) equations. The eight normalization
entries have weights \(0,0,1,-1,2,-2,1,-1\). The resulting equation
counts in weights \((-1,-2,1,2,0)\) are
\[
 (g_{-1}+3,\ g_{-2}+1,\ g_1+3,\ g_2+1,\ g_0+3).
\]
\src{src/Development/AutomorphicArithmetic/AutomorphicVarieties/UnitaryThreeGroups/BoxPresOf.lean}{149}{The rank-three framed ring}
and \src{src/Development/AutomorphicArithmetic/AutomorphicVarieties/UnitaryThreeGroups/MopBase.lean}{105}{the normalization entries}.

In rank four put \(u=(1,0)\), \(v=(-1,1)\), \(p=(0,1)\),
\(t=(1,1)\). The positive roots are \(u,v,p,t\). The zero
weight space has dimension two, and each of the eight root spaces has
dimension one. On weight \(j=(a,b)\), \(G\) acts by
\(\bar\alpha^a\bar\Psi^b\) and \(c\) by \((-1)^a\).
For \(j=0,\pm p\) this character differs from \(\bar\chi_\ell\)
at \(c\). For the six roots \(\pm u,\pm v,\pm t\), the specified
\(s_\ell\) has cyclotomic residue 1 and root-character residue
different from 1. Thus the twisted dual has no invariants and
\[
 o_j=g_j+2{\bf1}_{j=0}+{\bf1}_{j=p\ \mathrm{or}\ j=-p}.
\]
There is a separate arithmetic premise for \(g_j=0\) at the six
roots. It is the strict residual class vanishing for the characters
\(\bar\alpha^{\pm1},\bar\beta^{\pm1},\bar\theta^{\pm1}\) imposed by
the good datum. It kills \(\ker r_1\); the local invariants at every
\(v\in T\) vanish because the root character is nontrivial at \(s_v\).
The first dimension formula of the proof now gives \(g_j=0\).
\src{src/Development/AutomorphicArithmetic/AutomorphicVarieties/SymplecticGSp4/NumsZero.lean}{17}{The residual vanishing characters and their use in the count}.

\subsection{The rank-four local chart at the diagonal point}
\label{dp:chart}

Write \(\chi\) for the \(\ell\)-adic cyclotomic character.
The base characters in order are
\[
 (\theta,\Psi,\alpha,1),\qquad
 \Psi=\chi^e\psi,\quad\alpha=\chi\omega^{-1}\xi,\quad
 \beta=\Psi/\alpha,\quad\theta=\alpha\Psi .
\]
The multiplier is \(\theta\). The good datum has \(e\ge2\), finite
characters \(\psi,\xi,\omega\), and \(\xi\) unramified at \(T\).
At \(\ell\) it gives \(s_0\in I_\ell\cap\ker\psi\cap\ker\omega\)
with \(\chi(s_0)\) of infinite order. It also gives
\(\chi(s_\ell)=1\), \(\bar\omega=\bar\chi\), and
\(\bar\xi(s_\ell)\bar\Psi(s_\ell)^i\ne1\) for every integer \(i\).

\begin{proposition}[Local charts of the framed symplectic ring]\label{dp:local}
At the diagonal point, the completed characteristic-zero local ring
for the ordinary projector equations at \(\ell\) is a weighted power
series ring over \(E\). It has two coordinates of each weight
\(0,u,v,p,t\), and one of each weight \(-u,-v,-p,-t\).
For \(q\in S_b\) the corresponding ring for the condition
\(r(I_q\cap\ker\psi)=1\) has two weight-zero coordinates and one
of each of the eight root weights.
\end{proposition}

\begin{proof}
Use finite local Artin \(E\)-algebras \(A\) with residue \(E\),
given their finite-dimensional \(\ell\)-adic vector space topology.
At \(\ell\) write \(d_i=\psi_i(s_0)\) for the four diagonal values.
They are distinct, since their exponents on \(\chi(s_0)\) are
\(e+1,e,1,0\). The source equations specify
\(\det(X-r(s_0))=\prod_i(X-d_i)\). Hence
\[
 P_i=\prod_{k\ne i}\frac{r(s_0)-d_k}{d_i-d_k}
\]
are orthogonal projectors summing to one. Each image is a direct
summand of rank one by reduction to \(E\). The other equations say
\(P_jr(g)P_i=0\) for \(j>i\), and
\(P_i(r(h)-\psi_i(h))P_i=0\) for \(h\in I_\ell\). They specify
stability of the ordered flag and the prescribed inertial action on
its lines. The multiplier pairs lines \(i\) and \(5-i\) perfectly:
pairing any two eigenlines with \(d_i d_j\ne\theta(s_0)\) gives zero,
and nonsingularity of the entire pairing gives the assertion.
Rescaling eigenvectors lifting the original basis now produces a
symplectic basis reducing to that basis. Conversely an upper
triangular lift in such a basis with the given inertia characters
obeys the projector equations, since the partial sums of the
eigenlines of \(r(s_0)\) are its prescribed flag. The flag is
uniquely recovered by the projectors, including over nonreduced \(A\).

In the open cell of the symplectic flag variety at this flag there
is a unique lower unipotent matrix carrying the standard flag to a
given flag. This is Gaussian elimination with diagonal pivots units:
factor a near-identity change of symplectic basis into a lower
unipotent and an upper triangular factor. The involution
\(g\mapsto J^{-1}g^{-t}J\) preserves the two factors, so uniqueness
in \(\mathrm{GL}_4\) makes the factors symplectic. The lower
unipotent Lie algebra has one line in each negative root. Exponential
and logarithm are polynomial inverse maps here in characteristic
zero, since strictly triangular matrices are nilpotent. Thus the
moving-flag parameters give four smooth negative-weight directions.
It remains to consider a lift in the upper triangular group with
standard flag and fixed multiplier.

The diagonal deformations of \(\alpha\) and 1 are independent
unramified twists, determining the other two by the multiplier.
They lift across every small extension \(A'\to A\): use logarithm
on \(1+\mathfrak m_A\), lift the value at Frobenius additively, and
use \(\widehat{\mathbf Z}\to\mathbf Z_\ell\) to define the lifted
continuous additive character. Here and below a small extension has
kernel \(V\) killed by the maximal ideal of \(A'\). With lifted
diagonal fixed, the obstruction to lifting the upper-unipotent part
is in
\[
 H^2(G_{\mathbf Q_\ell},\mathfrak n\otimes_E V),\qquad
 \mathfrak n=E(\alpha)\oplus E(\beta)\oplus E(\Psi)\oplus E(\theta).
\]
Indeed a continuous vector-space section on the nilpotent entries,
in upper-unipotent exponential coordinates, gives a lifted cochain.
The multiplication error takes values in the additive kernel.
Conjugation on that kernel uses only the residual diagonal, so the
four root spaces split as written. A coboundary error can be
removed by subtracting a continuous one-cochain. Local Tate duality
gives \(H^2(G_{\mathbf Q_\ell},E(\gamma))^\vee
=H^0(G_{\mathbf Q_\ell},E(\chi/\gamma))\).
For \(\gamma=\alpha,\beta,\theta\) the residual value at \(s_\ell\)
differs from 1 by the displayed nontriviality with \(i=0,-1,1\),
respectively, while \(\chi(s_\ell)=1\). For \(\gamma=\Psi\) use
\(\chi(s_0)^{e-1}\ne1\). The four obstructions vanish. Every
surjection of finite local Artin \(E\)-algebras factors into small
extensions, giving formal smoothness.

The framed tangent space for a positive root \(\gamma\) consists of
all \(Z^1(G_{\mathbf Q_\ell},E(\gamma))\). By local Euler
characteristic and the preceding \(H^2=0\), its dimension is
\((1+h^0)+(1-h^0)=2\). The two unramified characters give two
zero-weight parameters. Together with the moving flag this gives
fourteen. The \(\ell\)-adic vector-space duality and Euler formula
used here follow from the finite-module theorems by applying them
to the finite quotients of a stable lattice and inverting \(\ell\).
For this passage the inverse systems of finite cohomology groups
are Mittag--Leffler, and every continuous cochain on the compact
powers of the Galois group takes values in a bounded lattice.

At \(q\in S_b\) the effective local group is \(G_q/(I_q\cap\ker\psi)\).
Its inertia quotient \(\Delta\) is finite since \(\ker\psi\) is
open. Its order may be divisible by \(\ell\). Over \(E\), averaging
over \(\Delta\) is exact. Consequently for each weight space \(W\)
of the symplectic adjoint, \(H^2\) of the effective group is zero
and \(H^1\) is the space of Frobenius coinvariants of \(W^\Delta\).
Indeed a stable lattice in \(W^\Delta\) reduces this last calculation
to cohomological dimension one of \(\widehat{\mathbf Z}\) on its
finite \(\ell\)-power quotients. Thus \(h^1=h^0\).
The small-extension multiplication-error argument with smooth
symplectic matrices gives formal smoothness, and its framed tangent
space has dimension \(h^1+\dim W-h^0=\dim W\) weight by weight.

To pass from this calculation to the completion, the integral local
ring \(R\) represents the actual matrix equations: both the
undivided projector expressions and the inertia-kernel conditions
are closed matrix equations in the unrestricted framed ring.
Write \(x:R\to O\) for the diagonal point and \(J=\ker x\).
Then \(\widehat{(R[1/\ell])_{J[1/\ell]}}\) represents precisely the
above Artin \(E\)-valued lifts. For such a lift, its matrix entries
minus their diagonal base values form a bounded subset of the
nilpotent ideal of \(A\) by compactness. Put them into a finite
\(O\)-lattice \(L\) in that ideal. If \(\mathfrak m_A^d=0\), the
order \(O+L+\cdots+L^{d-1}\) is complete local with residue \(k\)
and contains the entire representation. Its finite \(O\)-module
topology is its maximal-ideal topology, since its nilpotent ideal
is killed by the \(d\)-th power. This gives a map from \(R\)
reducing to \(x\), hence a local map on the indicated completion.
Conversely such a map kills \(J^d\), and \(R/J^d\) is a finite
\(O\)-module by the filtration with quotients \(J^i/J^{i+1}\).
It thus induces a continuous lift with bounded entries. Two
choices of order embed in a common one obtained by enlarging \(L\),
so the universal property also gives uniqueness.

Choose homogeneous lifts of a basis of \( (J/J^2)\otimes_O E\)
in \(J[1/\ell]\), using shifted matrix entries. They give a
surjection from a power series ring onto the completion. Formal
smoothness forces it to be an isomorphism: lift the inverse tangent
map successively to the power series ring modulo its maximal-ideal
powers. The resulting composite on the power series ring has
invertible linear term and is an automorphism by recursive
substitution. The coordinates and inverse are homogeneous in each
finite order, by uniqueness of the inverse. This also supplies
cotangent injectivity for the original integral ring after tensoring
with \(E\). The kernel of the map \(R\) to this completion is the
kernel of \(R\to(R[1/\ell])_{J[1/\ell]}\) by Krull intersection.
It is a finitely generated ideal, so a single \(y\in R\) with
\(x(y)\ne0\) annihilates it, by clearing the finitely many
denominators separately and taking their product.
\src{src/Development/AutomorphicArithmetic/AutomorphicVarieties/PadicHodgeTheory/SpFlag.lean}{309}{The inertia projector equations},
\src{src/Development/AutomorphicArithmetic/AutomorphicVarieties/PadicHodgeTheory/SpLiftL.lean}{32}{the four character checks}
and \src{src/Development/AutomorphicArithmetic/AutomorphicVarieties/SymplecticGSp4/LocBase.lean}{142}{the integral chart requirements}.
\end{proof}

\subsection{Cotangent faithfulness after the symplectic cuts}
\label{dp:cut}

\begin{lemma}\label{dp:cutlemma}
Let \(P\) be an augmented noetherian complete local \(O\)-algebra,
with augmentation ideal \(J\). Suppose a weighted chart
\(\phi:P\to E[[Y_1,\ldots,Y_n]]\) induces an isomorphism from
\((J/J^2)\otimes_O E\) to the cotangent at \(Y=0\).
Let \(u_1,\ldots,u_h\in J\) be homogeneous with independent images in
this rational cotangent. There is a chart of \(P'=P/(u_i)\) with
one fewer variable of the given weight for each cut, again inducing
a cotangent isomorphism. If \(P'\) is further divided by the kernel
of this latter chart, every \(O\)-valued point derivation at its
augmentation descends to the resulting domain.
\end{lemma}

\begin{proof}
Complete the linearly independent \(d\phi(u_i)\) to a basis in each
weight. Replace the first coordinates by the full series
\(\phi(u_i)\), keeping the complementary homogeneous linear
coordinates. The substitution has zero constant terms and invertible
linear part. Its inverse is obtained recursively in total degree.
Each finite degree respects the weights even when negative and zero
weights occur. Set the first \(h\) new coordinates to zero. On
cotangents this kills precisely the span of \(du_i\), and the
cotangent of \(P'\) is \(J/(J^2+(u_i))\). Exact scalar extension
therefore gives the asserted isomorphism. For \(a\) in the new
chart kernel its augmentation vanishes and its differential is zero
after inverting \(\ell\). Some \(0\ne c\in O\) then satisfies
\(ca\in(\ker x_{P'})^2\). Every integral point derivation kills this
square, hence kills \(a\) by torsion-freeness of \(O\). The chart
target is a domain.
\end{proof}

For the application take \(P=B[[t_1,\ldots,t_g]]\) from
Proposition~\ref{dp:relative}. The simultaneous local ring \(B\)
represents the product of the finitely many local functors. Its
rational completion is the power series ring on the union of the
local coordinates by Proposition~\ref{dp:local}. The relative
variables add their own free summands to the cotangent, so \(P\)
has the required chart. The global relations cut out \(U\) inside
this ambient formal spectrum. The next cuts are lifted to \(P\)
from the following entries of \(U\), in one-based indices:
\[
 {\cal P}=\big((1,2),(1,4),(3,2),(2,1),(4,1),(2,3)\big).
\]
At \(\ell\) take these entries of \(r(\phi_0)\), where
\(\phi_0=\lim_n s_\ell^{\ell^{n!}}\) is the element specified by the
datum. At every other \(v\in T\) take these entries of \(\alpha_v\).
All have zero augmentation. Use the six Lie algebra elements,
in the same order,
\[
 E_{12}-E_{34},\quad E_{14},\quad E_{32},\quad
 E_{21}-E_{43},\quad E_{41},\quad E_{23}.
\]
Each lies in \(\mathfrak {sp}_4(O)\) for the antidiagonal pairing with
signs \(1,1,-1,-1\), and has just one nonzero entry among \({\cal P}\).
For a selected element \(X\) conjugate \(r_0\) and all the local lifts
by \(1+\epsilon X\) over \(O[\epsilon]/(\epsilon^2)\), leaving all
frames equal to 1. This preserves the local equations by conjugation.
At \(\ell\) the only nonzero selected derivative is
\(\psi_j(\phi_0)-\psi_i(\phi_0)\). Indeed the residual diagonal at
\(\phi_0\) is the same as at \(s_\ell\), since the values belong to
\(k^\times\) and the residue degree \([k:\mathbf F_\ell]\) divides
\(n!\) eventually. Each pair has different first weight, so
\(\bar\xi(s_\ell)\bar\Psi(s_\ell)^m\ne1\) for all \(m\) makes
the residual entries distinct. At \(v\ne\ell\) keep the global
lift fixed, move only \(\alpha_v\) to \(1+\epsilon X\), and
conjugate the corresponding local diagonal lift by its inverse.
This tests one entry with derivative 1. These point derivations of
\(U\) pull back to \(P\), proving independence for any homogeneous
lifts of all the indicated cuts there.

Lemma~\ref{dp:cutlemma} now gives the post-cut chart and the
integral descent of derivations needed in
Section~\ref{dp:section}'s framed construction. Let \(b=|S_b|\).
In each odd root the cut removes \(b+1\) variables. The remaining
dimensions there are 1 for \(u,v,t\) and 0 for their negatives,
using the residual vanishings above. In weights \((-p,p,0)\)
the dimensions are
\[
 (1+b+g_{-p},\quad 2+b+g_p,\quad 2+2b+g_0).
\]
The remaining frame equations at each bad prime have weights
\(p,-p,0,0\), and the two normalization entries of
\(r(g_0)+\theta(g_0)r(g_0^{-1})\) have weights \(p,-p\).
Adding their counts to Proposition~\ref{dp:count} gives final lists
of lengths \(2+b+g_{-p},\,2+b+g_p,\,2+2b+g_0\), respectively.
They define the quotient of the resulting domain used by the model.
\src{src/Development/AutomorphicArithmetic/AutomorphicVarieties/SymplecticGSp4/ExModel.lean}{155}{The entries and six root vectors}
and \src{src/Development/AutomorphicArithmetic/AutomorphicVarieties/SymplecticGSp4/ExModel.lean}{425}{the normalized model}.

\subsection{The selected rank-three branch at a resonant character}
\label{dp:ordinary3}

Here \(G=G_{\mathbf Q_\ell}\), \(I\) is its inertia, and \(\eta\)
is an unramified character of order dividing two. Put \(H=\ker\eta\)
and \(I_0=I\cap H\cap\ker\tau\), where \(\tau\) has finite image.
The source specifies \(\Psi|_I=\chi^e\tau|_I\) with \(e\ge1\).
The matrix group is \(\mathrm{GL}_3\rtimes\{1,-1\}\), where
\(-1\) sends \(A\) to \(J A^{-t}J\) and \(J\) reverses the basis.
Put \((d_1,d_2,d_3)=(\Psi,1,\Psi^{-1})\). The base point of
the unrestricted local framed ring \(R\) is
\(r_0(g)=(\operatorname{diag}(d_i(g)),\eta(g))\), and its
augmentation is \(x\). For \(s\in I_0\) put \(A_i(s)=r(s)-d_i(s)\)
in the matrix block. The ordinary ideal \(I_{\rm ord}\) fixes the
characteristic polynomial on \(I_0\) to \(\prod_i(X-d_i(s))\)
and kills all entries of \(A_1(s)A_2(s')A_3(s'')\) for independent
\(s,s',s''\in I_0\). The definition then saturates by a specific
element \(\delta\), described below. Write
\[
 I_Y=(I_{\rm ord}:\delta^\infty),\qquad
 I_t=\bigcap_{\substack{\mathfrak p\text{ minimal over }I_Y\\
                          \mathfrak p\subseteq\ker x}}\mathfrak p .
\]
Taking the radical of \(I_Y\) before selecting these primes gives
the same intersection. The framed universal property makes \(x\)
the unique map with this specified matrix point.
\src{src/Development/AutomorphicArithmetic/AutomorphicVarieties/PadicHodgeTheory/IsInertiaInfiniteCyclotomic.lean}{97}{The ordinary equations, saturation and selected primes}.

\begin{proposition}\label{dp:rank3chart}
There is a unique selected prime. The completion of
\((R/I_t)[1/\ell]\) at \(x\) is a power series ring over \(E\)
with coordinate counts \(1,2,3,4,2\) in weights \(-2,-1,0,1,2\).
The kernel of the resulting chart from \(R\) is \(I_t\).
These statements hold also when a ratio of two adjacent diagonal
characters is cyclotomic.
\end{proposition}

\begin{proof}
As in Proposition~\ref{dp:local}, use continuous lifts over Artin
\(E\)-algebras, now of \(r_0\). Choose the specified
\(s_0\in I_0\) with \(\chi(s_0)\) of infinite order.
The eigenvalues \(d_i(s_0)\) are distinct. The characteristic
polynomial gives three rank-one projectors. Write \(V_1\) for the
first eigenline and \(V_2\) for the sum of the first two.
In eigen-coordinates, \(A_2(s_0)A_3(s_0)\) is a unit multiple of
the first projector, and \(A_1(s_0)A_2(s_0)\) is a unit multiple
of the third. The triple products therefore say that \(r(s)\) acts by
\(d_1(s)\) on \(V_1\), preserves \(V_2\), and acts on the last
quotient by \(d_3(s)\). Subtracting these two diagonal terms from
the specified trace gives \(d_2(s)\) on \(V_2/V_1\). Conversely
these three conditions imply every triple product, by applying its
factors from right to left along this flag.

The flag is unique by its ordered characters at \(s_0\). Conjugation
by \(g\in G\) preserves \(I_0\) and all the \(d_i\). When
\(\eta(g)=-1\) the induced action on flags first takes the dual flag
in reversed order via \(J\); its ordered quotient characters on
inertia are again \(d_1,d_2,d_3\) since \(d_{4-i}^{-1}=d_i\).
Uniqueness therefore gives stability under the full semidirect
action. By lower-unipotent change of basis the moving flag contributes
three free coordinates, two of weight \(-1\) and one of weight
\(-2\). Once the flag is standard all matrix blocks are upper
triangular, since inverse transpose across the antidiagonal preserves
the standard flag. Conjugation here and below is in the semidirect
group, so a change of basis also changes the block of an element
outside \(H\) by the prescribed involution.

The diagonal part relative to \(r_0\) is an unramified cocycle in
the nilpotent diagonal torus: on inertia the action on the torus
is trivial, and a character trivial on \(I_0\) factors there
through a finite group with no nonzero additive image in an
\(E\)-algebra. Take logarithms. The three coordinates \(b_i\) of
the value at a chosen Frobenius are free, even if \(\eta\ne1\).
Indeed a cocycle on \(\widehat{\mathbf Z}\) with Frobenius action
\(D=1\) or \(D(b_1,b_2,b_3)=(-b_3,-b_2,-b_1)\) has value
\(\sum_{k=0}^{n-1}D^kb\) at its \(n\)-th power. This extends
continuously using parity and the projection to \(\mathbf Z_\ell\).
Thus the base for the upper unipotent part is
\(\Lambda=E[[b_1,b_2,b_3]]\).

First divide the upper unipotent group by its central line \(E_{13}\).
The two adjacent entries form an additive module. For fixed diagonal
they are precisely continuous one-cocycles, with no quotient by
coboundaries. If \(H=G\), they give two independent rank-one
problems with respective parameters \(t=b_1-b_2,b_2-b_3\).
Their characters are \(\Psi\exp(t\nu)\), where the continuous
unramified additive character \(\nu\) is 1 on Frobenius.
If \([G:H]=2\), the two adjacent lines are exchanged by the
antidiagonal inverse transpose on their Lie algebra. The module
is induced from the first line on \(H\): this line and its translate
by a Frobenius form a basis, also for the deformed diagonal over
every Artin algebra. Its character on \(H\) is
\(\Psi|_H\exp(t\nu_H)\), with \(t=b_1-b_3\) and \(\nu_H\) equal
to 1 on Frobenius of \(H\). Indeed the value of the unramified
diagonal logarithm on the square of Frobenius is \(b+Db\).
This description by induction works functorially in the diagonal
parameters. In particular at \(t=0\) the indicated module is
isomorphic over the whole diagonal base to the constant module.

Here is the local equation for either one of these problems. Let
\(G_*=G\) in the split case and \(G_*=H\) in the induced case.
Denote the corresponding \(G\)-module at the origin by \(W\) and
put \(m=\dim_E Z^1(G,W)\). For its cocycle functor over \(\Lambda\),
a minimal presentation uses \(m\) variables and at most
\(\dim_E H^2(G,W)\) relations. To justify this version over \(E\),
the functor preserves fibre products of Artin algebras since its
objects are functions subject to equalities, and its tangent is
finite-dimensional by local Euler characteristic. Schlessinger's
prorepresentability criterion~\cite[Theorem 2.11]{schlessingerArtin}
therefore applies (the functor at the residue field is
one point). One can include the three diagonal
parameters in the functor and then work relative to their power
series ring. The small-extension argument in the proof of
Proposition~\ref{dp:relative} applies verbatim using continuous
linear sections of surjections of finite-dimensional \(E\)-spaces;
here the single obstruction is the coboundary of a lifted one-cochain.
The relation bound concerns the actual ideal modulo its maximal-ideal
multiple. Local Tate duality, with Shapiro's lemma in the induced
case, gives
\[
 \dim_E H^2(G,W)=
 {\bf1}_{\,\Psi|_{G_*}=\chi|_{G_*}} .
\]
If this is zero the cocycle algebra is a power series ring.

If the dimension is one, its possible single equation is, up to a
change of variables, exactly \(tY=0\). We verify its order of
vanishing. Choose the relative variables \(X_1,\ldots,X_m\) among
linear evaluations of the cocycle; they have weight one for scalar
multiplication of the cocycle. Modulo \(t\) the module is constant
up to the functorial basis isomorphism just proved, so every cocycle
is an arbitrary linear combination of a basis of \(Z^1(G,W)\) with
coefficients in the maximal ideal of the test algebra. Its completed
algebra over \(\Lambda/(t)\) is free in the \(m\) variables. The
chosen evaluations still give coordinates there by the inverse
function theorem. It follows that every equation is divisible by
\(t\). No relation has an \(X\)-independent term, since zero is
a cocycle for every diagonal.

Vary \(t\) to \(h\) and the cocycle to \(kz\) over
\(E[h,k]/(h^2,k^2,hk)\), choosing any linear diagonal section of
the parameter \(t\). The obstruction to lifting to
\(E[h,k]/(h^2,k^2)\), in the \(hk\) coefficient, maps to
\(\nu_{G_*}\smile [z_*]\) under the preceding cohomology
identification. In the split case this follows by expanding the
cocycle identity with character \(\chi(1+h\nu)\): its error on
\((g,g')\) is \(hk\,\nu(g)\chi(g)z(g')\) up to the sign of the
coboundary convention. In the induced case the cochain map to
\(G_*\) restricts and projects to the inducing line; it commutes
with reduction for the induced basis and gives the identical
calculation. Every class of \(H^1(G_*,E(\chi))\) occurs by
Shapiro. The local Tate pairing of this space with \(H^1(G_*,E)\)
is perfect and \(\nu_{G_*}\ne0\). Hence for some \(z\) the
obstruction is nonzero.

In particular there is one relation, with nonzero mixed term
\(t\sum_i c_i X_i\), some \(c_i\ne0\). It may be chosen
homogeneous of weight one. Indeed the ideal is stable under scalar
multiplication of all \(X_i\), which fixes \(b\). Its weight
components belong to the ideal by polynomial interpolation in each
finite truncation over the infinite field \(E\). All relations lie
in \(t(X_1,\ldots,X_m)\), so a relation with this nonzero
quadratic term generates the one-dimensional space of relations
modulo the maximal-ideal multiple. Its weight-one component still
generates. Write it \(t\sum_i a_i(b)X_i\); some \(a_i\) is a
unit. Use \(Y=\sum_i a_i(b)X_i\) as one of the relative
coordinates. This proves the claimed equation including its
multiplicity.

Finally restore the central root. Its Galois action on \(I_0\)
is \(\Psi^2=\chi^{2e}\), hence its \(H^2\) vanishes by local
duality since \(2e\ne1\). Multiplication errors for a lifted
central entry in a small extension are two-cocycles for this root
module, as above. The map to the first-layer functor is therefore
formally smooth with relative tangent its space of one-cocycles.
This adds free variables also over the singular first-layer ring:
choose relative cotangent coordinates and split the corresponding
power series surjection order by order by formal smoothness.
Equivalently the small-extension relation bound with zero obstruction
gives zero relative relations. The central coordinate has weight two.

We compute the effect of the prescribed saturation in these
coordinates. The source chooses \(f\in H\cap\ker\tau\) mapping to
a positive power of Frobenius. For \(i=1,2,3\) put
\(c_i=\prod_{k\ne i}(d_i(s_0)-d_k(s_0))\in E^\times\).
Its trace expression is \(c_i\,\mathrm{tr}(P_i r(f))\), where
\(P_i\) is the corresponding eigenprojector of \(r(s_0)\). In an
upper triangular basis this equals \(c_i\kappa_i(f)\), with
\(\kappa_i\) the \(i\)-th diagonal character on \(H\). Indeed each
\(P_i\) is upper triangular with only its \(i\)-th diagonal entry
equal to 1. Thus \(\delta\) is, up to a unit in the characteristic-zero
completion, precisely
\[
  \prod_{i=1}^2\bigl(\kappa_i(f)/\kappa_{i+1}(f)-\chi(f)\bigr).
\]
For a resonant adjacent block the corresponding factor is
\(\chi(f)(\exp(a t)-1)\), with \(a=\nu_{G_*}(f)>0\).
It is a unit times \(t\). This formula also covers the two equal
ratios on \(H\) in the induced case. In the quotient by all the
coordinates \(Y\) belonging to resonant blocks, the ring is a
power series domain and the image of \(\delta\) is nonzero:
each adjacent ratio still varies as
\(\Psi(f)\exp(a t)\) for a nonzero diagonal parameter \(t\).
Since \(tY=0\) and \(t\) divides \(\delta\) in each resonant
block, the \(\delta\)-power torsion is exactly the ideal of these
\(Y\)'s. This establishes the calculation even if \(\delta(x)=0\).

For any \(G\)-module \(W\) in an adjacent block the local Euler
formula over \(\mathbf Q_\ell\) gives
\(\dim Z^1(G,W)=2\dim W+\dim H^2(G,W)\).
Saturation removes exactly one variable for each of the latter
dimensions. The adjacent roots thus leave four weight-one
coordinates in all; the central root adds two of weight two.
The three unramified diagonal coordinates and the three lower flag
coordinates give the other asserted counts.

Localization at the characteristic-zero point and completion are
flat and commute with \((I_{\rm ord}:\delta^n)\): take the kernel
of multiplication by \(\delta^n\) on the quotient. The chain
stabilizes by noetherianity. It follows that the completion of the
localization modulo \(I_Y\) is the power series domain just computed.
Faithful flatness makes the localized quotient a domain.
Consequently there is exactly one minimal prime over \(I_Y\)
contained in \(\ker x\), the contraction of the zero ideal of this
local quotient. Its contraction to \(R\) is also the kernel of the
completed chart. This proves the assertion about the selected
integral ideal \(I_t\).
\src{src/Development/AutomorphicArithmetic/AutomorphicVarieties/PadicHodgeTheory/RtChart.lean}{41}{The corresponding twelve-coordinate conclusion}.
\end{proof}

The preceding computation is a statement about the ordinary matrix
ideal in the completed local ring. Application to the integral
presentation uses a weight grading on \(R/I_t\). The unrestricted
matrix ring has homogeneous coordinates given by shifted entries.
The ordinary equations are weight-homogeneous, \(\delta\) has
weight zero, and the map to the completed chart preserves weights.
Its kernel \(I_t\) is thus closed under projection onto each weight.
Conjugation by a matrix congruent to 1 also preserves the selected
branch over arbitrary integral test algebras. Indeed conjugate the
universal \(R/I_t\)-valued family by the generic matrix \(1+X\)
in \((R/I_t)[[X_{ij}]]\). Embed this ring coefficientwise into
\(E[[Y_1,\ldots,Y_{12},X_{ij}]]\) by the chart. The ordinary
equations are conjugation invariant; \(\delta\) is unchanged because
its defining traces use elements of \(H\), and it is nonzero in this
domain. Hence the conjugate annihilates \(I_Y\). For \(a\in I_t\)
some \(s\) with \(x(s)\ne0\) satisfies \(sa\in I_Y\), by the
calculation of the localized prime. The conjugate image of \(s\) has
constant coefficient \(x(s)\), so \(a\) also maps to zero.
Injectivity then gives the same conclusion before passage to \(E\).
Substituting any matrix of elements of the maximal ideal for \(X\)
proves the integral assertion. This supplies the local presentation required by
Proposition~\ref{dp:relative} in rank three. Identifying potentially
crystalline points of the required type on the entire model is
a further arithmetic statement.

In rank four the local and relative arguments use the good-datum
hypotheses, including the separate residual vanishings.
In both ranks, finiteness of the integral invariant quotient,
exclusion of extraneous one-sided points and rigidity of the
nontrivial two-sided points still require arithmetic control of
representations away from the diagonal base point. The local charts
and equation counts alone yield none of these global-point assertions.

\subsection{The special fibre of the invariant algebra}\label{dp:finitecurves}

There is a separate arithmetic test for the finiteness hypothesis in
Theorem~\ref{newcoh:orbitdegree}. In rank three it leaves only
strongly irreducible three-dimensional specializations on \(G_K\).
In rank four it leaves an absolutely irreducible four-dimensional
specialization or an even absolutely irreducible two-dimensional
constituent. The fields of coefficients in these statements have
characteristic \(\ell\).

\begin{lemma}[Detecting a nonfinite invariant algebra]\label{dp:detect}
Let \(O\) be a complete DVR with finite residue field \(k\) and
uniformizer \(\pi\). Let \(S\) be a quotient of a finite formal
power series algebra over \(O\) by a closed ideal stable under
projection onto the weight components for integral vector weights.
Write \(C\subset S\) for the image of the weight-zero power series.
Assume the variables are in the maximal ideal. Then \(C\) is finite
over \(O\) if and only if every continuous local map
\[
               S\longrightarrow k'[[t]],\qquad \pi\longmapsto 0,
\]
for finite extensions \(k'/k\), takes \(C\) into the constants.
\end{lemma}
\begin{proof}
The nonzero minimal exponent vectors of weight zero are finite
by Dickson's lemma, and they generate the monoid of all balanced
exponents by subtraction. Denote the corresponding monomials by
\(m_j\). Substituting these into a formal power series over \(O\)
surjects onto \(C\): fix one factorization for each exponent and
group the terms accordingly. The bounded total degrees of the
\(m_j\) ensure convergence in both descriptions. Thus \(C\) is
complete noetherian local, with maximal ideal \((\pi,m_j)\).
Weight-zero projection gives \(C\cap\pi S=\pi C\).
If \(C/\pi C\) is artinian, a power of its maximal ideal is zero.
Consequently the \(\pi\)-adic and maximal-ideal topologies on \(C\)
are equivalent; lifting a \(k\)-basis of \(C/\pi C\) and expanding
successively in powers of \(\pi\) gives finitely many \(O\)-module
generators. Conversely a finite algebra has artinian special fibre,
whose image in the domain \(k'[[t]]\) is a finite field.

Otherwise some \(m_j\) is not nilpotent modulo \(\pi C\). By the
projection identity it survives on a prime quotient of \(S/\pi S\).
Cut this complete local domain to dimension one while preserving
the image \(b\ne0\) of \(m_j\). At each step of dimension \(>1\)
choose an element of the maximal ideal outside the finitely many
height-one primes dividing \(b\), and take a minimal prime over
that nonzero element. This prime has height one, avoids \(b\), and
is strictly below the maximal ideal. The resulting one-dimensional
complete domain has finite normalization since it is excellent.
This normalization is local by completeness and the domain property,
and is an equicharacteristic complete DVR with finite residue
extension. Its coefficient field is the unique lift of this finite
field, so a uniformizer identifies it with \(k'[[t]]\) over \(k\).
Here \(b\) has strictly positive valuation. This proves the criterion.
\end{proof}

In the two presentations one may use homogeneous coordinates taken
from matrix entries minus their residual diagonal lifts. Indeed these
entries generate the representing ring topologically, by the framed
universal property, and homogeneous cotangent generators suffice by
successive approximation. Equivalently, on the post-cut rank-four
model the entries generate \(S\) in the sense that they distinguish
its maps to complete local algebras. They then span its relative
cotangent space by testing the dual numbers of the residue field.
For these coordinates \(C\) is the closure of the algebra generated
by weight-zero products of entries. Closure does not change module
finiteness: a finite \(O\)-submodule is compact, hence closed, and
an \(O\)-submodule of a finite \(O\)-module is finite. Thus
Lemma~\ref{dp:detect} applies to the actual weight-zero entry algebra.
The image of the invariant part \(B\) of the presentation in the
post-cut ring is \(B/(I\cap B)\). In rank four one first takes
weight zero for the first coordinate and subsequently for the second;
this is simultaneous weight \((0,0)\).
\src{src/Development/AutomorphicArithmetic/AutomorphicVarieties/CurveDVRSpecialization/FinCrit.lean}{122}{Curve detection in rank three};
\src{src/Development/AutomorphicArithmetic/AutomorphicVarieties/SymplecticGSp4/FgMapMk.lean}{24}{From the rank-four weight-zero entry algebra to the invariant quotient}.

Here are two arithmetic facts about such curves. We may enlarge
\(k'((t))\) by a finite extension and normalize its valuation ring
to define constituent spaces, their flags and stable lattices.
Take intersections and saturated quotients of the original lattice
to get stable lattices on subquotients. Their residual constituents
together give the residual semisimplification of the original
representation, by reduction of these exact lattice sequences.

First, a continuous character of a number-field Galois group into
the units of an equal-characteristic DVR with finite residue field
is constant if its inertia characters at every finite place have
finite image, with ramification at only finitely many places.
Indeed a prime-to-\(\ell\) root of unity is the unique lift of its
residue and the principal units have no torsion in characteristic
\(\ell\). Dividing by the lift of the residual character leaves an
everywhere unramified character into principal units. It factors
through the finite narrow class group by global class field theory,
and is therefore trivial. This argument applies on every finite
number-field extension as well. If a matrix representation becomes
triangular there, and its inertial eigencharacters have finite
image, all its diagonal characters there are constant. On the full
group every eigenvalue is then a root of unity, by raising a group
element to a positive power in the subgroup. Its characteristic
polynomial is the constant lift of its reduction, even if the
subgroup and power depend on the curve.

\begin{lemma}[Odd two-dimensional subquotients on a curve]\label{dp:pan}
Let \(\ell>2\). Suppose a continuous two-dimensional representation
of \(G_{\mathbb Q}\) over a finite extension of \(k'((t))\) has a
stable lattice, is unramified outside finitely many primes, has
finite inertial eigencharacters everywhere, and is triangular on
\(G_{\mathbb Q_\ell}\) over an algebraic closure. Suppose its
residual semisimplification is a sum of two characters of opposite
parity. Then its characteristic polynomials agree with the constant
lifts of their reductions; in particular it is geometrically
reducible.
\end{lemma}
\begin{proof}
After a finite coefficient extension twist by the constant inverse
of the first residual character. Write \(1,\bar\chi\) for the
resulting residual characters. By the preceding character argument
the determinant is the constant lift of \(\bar\chi\). Use its
Teichm\"uller lift \(\chi\) to \(W(k')\), enlarging \(k'\) to the
residue field of the chosen extension. It is an odd finite-order
character and is de Rham locally at \(\ell\). We use the number
field \(\mathbf Q\), an abelian totally real field unramified at
\(\ell\), so the required extension of the residual determinant
to \(G_{\mathbf Q}\) is already given. Choose a finite set \(S\)
of finite primes containing \(\ell\) and every ramified prime
of the twisted representation. Pan's ordinary
pseudodeformation ring for \(G_{\mathbb Q,S}\) at \(1+\bar\chi\),
with determinant \(\chi\), is finite over
\(\Lambda=W(k')[[\mathbb Z_\ell^\times(\ell)]]\), where the group
in brackets is the pro-\(\ell\) part of the local units.
Use the distinguished theorem when
\(\bar\chi|_{G_{\mathbb Q_\ell}}\ne1\). If its restriction is
trivial, use the character-labelled theorem and the actual
subcharacter \(\lambda_1\) of the local triangular representation.
Both local diagonal characters are integral units, since the
inverses are integral too. In a triangular basis
\((r(\sigma)-\lambda_1(\sigma))(r(\tau)-\lambda_2(\tau))=0\)
for local \(\sigma,\tau\); taking the trace after multiplying by
any \(r(g)\) gives the extra ordinary pseudocharacter relation.
The coefficient homomorphism from \(W(k')\) to the normalized
curve DVR is reduction to \(k'\) followed by its constant inclusion.
Thus the trace and, in the second case, the chosen character
define the asserted map of universal rings on passing to the
inverse limit of the Artin quotients. The finiteness assertions
used here hold for every odd prime, including three.
The map on \(\Lambda\) comes from an inertia
character which has finite image in the residue-field units, so
it kills the augmentation of the pro-\(\ell\) group and also
\(\ell\). The image of the finite algebra in this DVR is a finite
field. The trace and determinant therefore equal their residual
values as constants. The characteristic-polynomial criterion for
semisimplification (Brauer--Nesbitt) finishes the proof.
\cite[Theorems~5.1.1(1) and~6.1.1(1) in the published
version, Theorems~5.1.2(1) and~6.1.2(1) on pp.~71 and~95--96
of the author manuscript]{panFM}
\end{proof}

We explain how traces control the marked entries required here. Put
\(\mathcal E=\operatorname{span}\{r(g)\}\) over the algebraic
closure of the fraction field of the curve, using \(G_K\) in rank
three. It is an algebra. It contains the individual coordinate
projectors \(E_{ii}\) in both models. In rank three the marked
difference \(Z=r(g_1)-Jr(g_1)^tJ\) belongs to this algebra since
the second matrix is \(r(cg_1^{-1}c)\). The normalization kills
entries \(01,10\), and antisymmetry under reverse transpose makes
\(Z=\operatorname{diag}(z,0,-z)\). The reduction of \(z\) is
\(\bar\Psi(g_1)-\bar\Psi(g_1)^{-1}\ne0\). Lagrange interpolation
gives all three projectors. In rank four the marked block matrix
\(r(\phi)\) has coprime characteristic polynomials on the blocks
\(\{0,2\}\) and \(\{1,3\}\), by the distinct residual values
between blocks. Polynomial interpolation gives the first block
projector. In that block \(r(g_0)+\theta(g_0)r(g_0)^{-1}\) is
diagonal by normalization, with diagonal difference reducing to
\((\bar\alpha(g_0)-1)(\bar\Psi(g_0)-1)\ne0\). We obtain
\(E_{00},E_{22}\) and then their symplectic adjoints \(E_{33},E_{11}\).
\src{src/Development/AutomorphicArithmetic/AutomorphicVarieties/SymplecticGSp4/EiiUnit.lean}{169}{The four marked projectors}.

If the semisimplification is a sum of characters, a full invariant
flag is consequently a coordinate flag: every stable subspace is
preserved by all \(E_{ii}\). Order the coordinates by this flag,
assigning numbers \(b_i\) so \(b_i>b_j\) for every nonzero
off-diagonal entry in position \(ij\). The reverse position
\((j^*,i^*)\), with \(i^*=n-1-i\), also occurs. This follows by
reverse transpose in rank three and symplectic adjunction in rank
four, both preserving \(\mathcal E\). Therefore \(h_i=b_i-b_{i^*}\)
is strictly decreasing on these directed entries. The functional
\(ij\mapsto h_i-h_j\) factors through their weights. For rank three
this follows from \((h_0,h_1,h_2)=(h_0,0,-h_0)\). For rank four it
is the linear form \((h_0-h_1)\lambda+(h_0-h_2)\mu\) on differences
of the coordinates \((\lambda_i,\mu_i)=(1,1),(0,1),(1,0),(0,0)\).
Thus a nonzero balanced monomial uses diagonal entries only. These
are the diagonal characters on the flag and are constant by class
field theory. In rank three all negative-sign elements have the
matrix of a positive-sign element, by multiplication by \(c\).
Conversely, an absolutely irreducible constituent of dimension \(>1\)
forces a nonconstant characteristic-polynomial coefficient if the
residual semisimplification consists of characters. Such coefficients
are themselves sums of weight-zero monomials.

\begin{proposition}[Remaining characteristic-\(\ell\) curves]\label{dp:curvebranches}
For finiteness of the invariant quotient in rank three it suffices
to exclude curves on which \(r|_{G_K}\) is strongly absolutely
irreducible. In rank four it suffices to exclude curves on which
\(r\) is absolutely irreducible, and those with an absolutely
irreducible two-dimensional constituent on which complex conjugation
is scalar.
\end{proposition}
\begin{proof}
The needed local flags specialize even where the marked inertia
eigenvalues coalesce modulo \(\pi\). In rank three take the closed
incidence scheme in the projective full flag variety requiring
stability under \(H\) and the ordered characters \(d_i\) on \(I_0\).
It has a point over the fraction field of the chart constructed
in Proposition~\ref{dp:rank3chart}. Its proper image contains the
generic point of \(R/I_t\) by the kernel assertion there, so every
geometric fibre has a flag. In rank four use the symplectic flag
variety, stability under the full decomposition group and the four
prescribed inertia characters. The condition \(\mathrm{OnComp}\)
gives a domain carrying the local equations and an augmentation to
\(O\) at the diagonal point, through which the test factors. Thus
the differences \(d_i(s_0)-d_j(s_0)\) are nonzero in its fraction
field. There the projectors give a flag by the calculation in
Proposition~\ref{dp:local}. Properness specializes it to the given
geometric point. The closed
conditions in both assertions are intersections of the closed
stability and action equations, so no finite presentation of the
inertia group is needed. Away from \(\ell\) the fixed residual
or finite character kernel bounds inertia. At \(\ell\) the
inertia characters on these flags have finite image in equal
characteristic, also on the full inertia in rank three since
\([I:I_0]\) is finite. These properties pass to subquotients.

In rank three the algebra \(\mathcal E\) contains \(E_{ij}\)
precisely when some matrix entry in that position is nonzero, by
sandwiching between \(E_{ii}\) and \(E_{jj}\). Its entry relation
is thus transitive. Its equivalence classes for mutual
reachability are the simple diagonal blocks after the classes
are topologically ordered, each supporting a full matrix algebra.
Reverse transpose permutes the classes and reverses the order.
The only possible two-element class is \(\{0,2\}\): its reverse
must be the unique class of size two. The middle singleton is
incomparable to that class, since both classes are fixed by
reversal. Thus a two-dimensional constituent forces the actual
outer block decomposition. The middle character and the determinant
on \(G_K\) have finite inertia and trivial reduction and hence are
one by the character argument. The outer block \(U\) has determinant
one. It extends to \(W\) on \(G_{\mathbb Q}\) with
\(W(c)=D=\operatorname{diag}(-1,1)\), by the matrix identity
\(J_2 U(h)^{-t}J_2=D U(h)D^{-1}\) for \(h\in G_K\), with
\(J_2\) the two-by-two antidiagonal unit matrix. The marked
outer coordinates retain the original integral lattice, which
is also preserved by \(D\). The formula on the two open cosets
is continuous, and \(D^2=1\) gives the group law
as in Proposition~\ref{spreal:end}.
Its residual characters are \(\eta\bar\Psi,\bar\Psi^{-1}\).

On \(I_\ell\subset G_K\) the two inertia constituents of \(W\)
are finite and hence constant, equal to \(\bar\Psi,\bar\Psi^{-1}\)
by reduction. They are distinct since \(\psi_\ell^2=1\) and
\(\ell-1\nmid2e\). The restriction to inertia has a common eigenline
by the local flag. The eigenspace for its character is exactly one
line by distinctness. Its character is invariant under the full
decomposition group because it is the restriction of one of the
global residual characters. The line is therefore stable under
that group. Away from \(\ell\) the extension adds only the finite
quadratic character. Lemma~\ref{dp:pan} now excludes an absolutely
irreducible outer block. If an irreducible three-dimensional
representation becomes reducible over a finite number-field
extension, pass to its normal core. Its restriction there is
semisimple: the sum of the simple submodules is nonzero and is
stable under the full group. The isotypic components are permuted
transitively, with equal multiplicities. Primality of dimension
three makes every simple constituent on this subgroup a character.
The preceding potential triangularity argument gives constant
characteristic polynomials, contradicting absolute irreducibility.

In rank four any reducible symplectic similitude representation
has simple constituents of dimension at most two. A stable line
lies in its stable three-dimensional perpendicular and a stable
hyperplane gives such a line by duality. A stable two-plane gives
two pieces of dimension two directly. Every two-dimensional
constituent inherits a local flag by intersection and quotient,
and its residual semisimplification is a pair from
\((\bar\theta,\bar\Psi,\bar\alpha,1)\). Lemma~\ref{dp:pan} excludes
all opposite-parity pairs. The equal-parity pairs are precisely
\(\{\bar\theta,\bar\alpha\}\) and \(\{\bar\Psi,1\}\), with parity
read on the two coordinate pairs \(\{0,2\}\), \(\{1,3\}\).
The preceding marked-entry argument and Lemma~\ref{dp:detect}
complete both reductions.
\end{proof}

The rank-four source separates this issue from its characteristic-zero
purity and automorphy assertions. Its first deduction of finiteness
uses the positive-characteristic exclusion \(\mathrm{Con}(D)\), whose
test includes the nonsplit extension with irreducible even block on
indices \(1,3\). In the tower the ideal \(J\) is required to satisfy
\(\mathfrak m_{X.B}^n\subset J+\mathfrak m_{O'}X.B+
(\psi_{v,i}(s)-1:s\in I_v)\). This gives finiteness on the
fixed-inertia special fibre \emph{of the quotient by \(J\)}.
Using it for a given curve also requires showing that the curve
annihilates \(J\) on the tower. Its localization clause for
irreducible, locally distinguished curves already assumes this
annihilation. The characteristic-zero seed and purity clauses
alone impose no equation on a purely characteristic-\(\ell\)
component. In rank three the source uses a common finite-index
subgroup for constancy of all the special-fibre characteristic
polynomials. The curve criterion shows why it is enough to
settle the remaining representations of Proposition~\ref{dp:curvebranches}
one curve at a time.
\src{src/Development/AutomorphicArithmetic/AutomorphicVarieties/CurveDVRSpecialization/GfResZpow.lean}{179}{Rank-three finiteness from special-fibre characteristic polynomials};
\src{src/Development/AutomorphicArithmetic/AutomorphicVarieties/ModularityArgument/ThresholdArguments/DenseThresholdCases/VWinPm.lean}{75}{The characteristic used in each of the three conclusions};
\src{src/Development/AutomorphicArithmetic/AutomorphicVarieties/ModularityArgument/ThresholdArguments/DenseThresholdCases/PotentialModularityTower.lean}{38}{Finiteness and localization conditions on the tower ideal}.
